\problem{1}
\textbf{[Отбор 2025, 8.1]}
Даны положительные числа $x, y$ такие, что $x+y=1$. Докажите неравенство:

$$
\frac{x+1}{y}+\frac{y+1}{x} \geq 6 .
$$

\problem{2}
\textbf{[Отбор 2024, 8.3]}
Чебурашка и Гена вместе отправились из магазина домой. Через 15 минут, когда они прошли четверть пути до дома, Чебурашка заметил, что забыл апельсины в магазине и побежал за ними. Забрав апельсины, он так спешил домой, что прибежал на 30 минут раньше Гены, который шёл со своей постоянной скоростью. Во сколько раз скорость бега Чебурашки больше скорости ходьбы Гены, если всю дорогу после расставания с Геной Чебурашка бежал с постоянной скоростью?

\problem{3}
\textbf{[Отбор 2023, 8.4]}
Пусть $n$ --- натуральное число. Каждое из чисел $1,2,3, \ldots, 100$ окрашено в один из $n$ цветов таким образом, что любые два различных числа, сумма которых делится на 4 , окрашены в разные цвета. Найдите наименьшее возможное значение $n$, при котором существует такая раскраска.

\problem{4}
\textbf{[Отбор 2025, 8.3]}
У Даши и Эмилии есть стопка из 2025 карт, и они играют в следующую игpy: ходят по очереди, и каждый игрок забирает сверху количество карт, которое должно быть некоторой степенью двойки, т.е. $1,2,4,8, \ldots$ Выигрывает игрок, который заберёт последнюю карту. Даша начинает игру. Верно ли, что для Эмилии существует стратегия, гарантирующая ей победу, независимо от того, как играет Даша?

\problem{5}
\textbf{[Отбор 2020, 8.6]}
Для скольких натуральных чисел $n$ сумма $\left[\frac{2019}{n}\right]+\left[\frac{2020}{n}\right]+\left[\frac{2021}{n}\right]$ не делится на 3? ($[x]$ --- это наибольшее целое число, не превосходящее $x$.)

\problem{6}
\textbf{[Отбор 2021, 8.6]}
Две криптовалюты были созданы 22 мая 2021 года. Каждая стоит сначала по 1 рублю за монету. Стоимость первой валюты увеличивается на 1 рубль в день. А курс второй валюты каждый день вычисляется как отношение суммы курсов двух валют за прошлый день к произведению этих же двух значений. Найдите курс второй валюты через 3650 дней.